% ============================================================================
% TABLA 2.6: CUADRO COMPARATIVO DE INVESTIGACIONES PREVIAS
% Generado por: Poseidón 🔱 (12 Nov 2025, 21:15 hrs)
% Selección: 17 artículos representativos (2018-2025)
% Criterios: Diversidad metodológica + LOUO destacado + Alcance amplio
% ============================================================================

\begin{table}[htbp]
    \centering
    \caption{Cuadro comparativo de investigaciones previas sobre clasificación de comportamiento sedentario y actividad física con wearables}
    \label{tab:comparativa_investigaciones}
    \footnotesize
    \begin{tabular}{@{}p{1.8cm}p{2.2cm}p{1.5cm}p{1.3cm}p{0.8cm}p{1.2cm}p{1.8cm}p{1.5cm}@{}}
        \toprule
        \textbf{Autor (Año)} & 
        \textbf{Objetivo} & 
        \textbf{Técnica} & 
        \textbf{Datos/ Wearable} & 
        \textbf{N} & 
        \textbf{Validación} & 
        \textbf{Métrica Principal} & 
        \textbf{Resultado} \\
        \midrule
        
        % ========== FUZZY LOGIC ==========
        
        Capitoli et al. (2025) \cite{Capitoli2025FuzzyXAI} & 
        Diagnóstico clínico interpretable con lógica difusa probabilística & 
        Fuzzy Logic (interpretable) & 
        Datos clínicos hospitalarios & 
        Clínico & 
        Validación clínica & 
        Interpretabilidad 100\% & 
        Clínicos controlan proceso diagnóstico. Fuzzy = XAI \\
        
        \midrule
        
        Marashi-Hosseini et al. (2023) \cite{MarashiHosseini2023Dietary} & 
        Sistema decisión dietética para pacientes múltiples condiciones crónicas & 
        Fuzzy Logic Mamdani (1144 reglas) & 
        Datos dietéticos + clínicos & 
        100 & 
        Validación vs nutricionistas & 
        Precisión & 
        97\% precisión. Sistema experto fuzzy riguroso \\
        
        \midrule
        
        Gonçalves et al. (2021) \cite{Goncalves2021} & 
        Reconocimiento actividad humana con clustering + fuzzy & 
        K-Means ($k=2$) $\rightarrow$ FIS Mamdani & 
        Sensores inerciales (IMU) & 
        --- & 
        Validación estabilidad & 
        Estabilidad clasificación & 
        ÚNICO precedente clustering $\rightarrow$ fuzzy supervisado \\
        
        \midrule
        
        % ========== DEEP LEARNING ==========
        
        Khan et al. (2024) \cite{Khan2024Wearable} & 
        Reconocimiento actividad física con sensores inerciales + deep learning & 
        Deep Learning (CNNs + LSTMs) & 
        Sensores inerciales wearable & 
        --- & 
        K-fold CV & 
        Accuracy & 
        Alta precisión. Requiere datos masivos + no interpretable \\
        
        \midrule
        
        Bassani et al. (2025) \cite{Bassani2025DLHAR} & 
        Deep Learning para reconocimiento actividad física en tareas manuales & 
        Deep Learning (varios modelos) & 
        Acelerómetros & 
        --- & 
        70/30 vs LOSO & 
        Accuracy & 
        \textbf{70/30: 95.7\%} $\rightarrow$ \textbf{LOSO: 90.3\%}. Generalización robusta requiere LOSO \\
        
        \midrule
        
        % ========== MACHINE LEARNING CLÁSICO ==========
        
        Fuller et al. (2021) \cite{Fuller2021Predicting} & 
        Clasificación lying, sitting, walking, running con Apple Watch + Fitbit & 
        Random Forest & 
        Apple Watch + Fitbit (acelerómetro) & 
        33 & 
        LOOCV & 
        Accuracy & 
        Accuracy $>$90\% para posturas. BYOD validado \\
        
        \midrule
        
        Mathew et al. (2024) \cite{Mathew2024LOSO} & 
        Medición uso funcional mano en parálisis cerebral unilateral con ML & 
        Machine Learning (varios modelos) & 
        Acelerómetros (muñeca) & 
        22 & 
        Personalizado vs LOSO & 
        F1-Score & 
        \textbf{Personalizado: F1=0.896} $\rightarrow$ \textbf{LOSO: F1=0.584}. Evidencia sobreajuste $N$ pequeño \\
        
        \midrule
        
        Rehman et al. (2024) \cite{Rehman2024LOSO} & 
        Estudio comparativo validación + extracción características en HAR wearables & 
        Random Forest & 
        Sensores inerciales wearables & 
        --- & 
        LOSO vs k-fold & 
        Accuracy & 
        \textbf{LOSO: 76\%} vs \textbf{k-fold: 89\%}. Data leakage justifica LOUO \\
        
        \midrule
        
        % ========== MÉTODOS ESTADÍSTICOS ==========
        
        Salim et al. (2024) \cite{Salim2024STEPHEN} & 
        Detección tiempo sedentario + bouts con Hidden Semi-Markov Model & 
        Hidden Semi-Markov Model & 
        Wearables consumo (muñeca) & 
        --- & 
        Validación vs ground truth & 
        Sensibilidad/ Especificidad & 
        Detección precisa sedentary bouts. Método estadístico avanzado \\
        
        \midrule
        
        Migueles et al. (2022) \cite{Migueles2022GRANADA} & 
        Consenso análisis acelerómetro para comportamientos físicos en epidemiología & 
        Consenso metodológico & 
        Acelerómetros (varios) & 
        --- & 
        Revisión sistemática & 
        Guías metodológicas & 
        Estándar internacional para análisis PA/sedentarismo \\
        
        \midrule
        
        % ========== VALIDACIÓN WEARABLES ==========
        
        Lyons et al. (2024) \cite{Lyons2024StandHour} & 
        Definición umbrales actividad que activan ``Stand Hour'' en Apple Watch & 
        Estudio observacional & 
        Apple Watch (energía activa + pasos) & 
        --- & 
        Análisis umbral & 
        Umbrales Stand Hour & 
        Deconstruye proxy sedentarismo comercial. Umbrales caja negra identificados \\
        
        \midrule
        
        O'Grady et al. (2024) \cite{OGrady2024AppleWatch} & 
        Validación HRV + RHR en Apple Watch Series 9 y Ultra 2 & 
        Estudio validación & 
        Apple Watch vs Polar H10 (gold standard) & 
        39 & 
        Validación concurrente & 
        MAPE & 
        \textbf{HRV MAPE=28.88\%}, \textbf{RHR MAPE=5.91\%}. Limitaciones sensor PPG \\
        
        \midrule
        
        Ji et al. (2023) \cite{Ji2023Scratch} & 
        Evaluación scratch nocturno con actigrafía en pacientes dermatitis atópica & 
        Machine Learning & 
        Actigrafía (ActiGraph) & 
        --- & 
        \textbf{LOSO evaluation} & 
        Sensibilidad/ Especificidad & 
        \textbf{Nature npj Digital Medicine}. Precedente metodológico clínico LOSO \\
        
        \midrule
        
        % ========== XAI + INTERPRETABILIDAD ==========
        
        Bienefeld et al. (2023) \cite{Bienefeld2023XAI} & 
        Resolver dilema XAI: brecha necesidades clínicos vs objetivos desarrolladores & 
        Estudio cualitativo & 
        Entrevistas + cuestionarios & 
        112 & 
        Análisis cualitativo & 
        Temas identificados & 
        Clínicos requieren transparencia global. Desarrolladores enfocan local \\
        
        \midrule
        
        Deng et al. (2023) \cite{Deng2023LharJBHI} & 
        HAR ligero en wearables con knowledge distillation & 
        Deep Learning (ligero) & 
        Sensores wearables & 
        --- & 
        Validación experimental & 
        Accuracy + eficiencia & 
        \textbf{IEEE JBHI} (revista objetivo). HAR eficiente wearables \\
        
        \midrule
        
        % ========== FISIOLOGÍA + HRV ==========
        
        Godkin et al. (2025) \cite{Godkin2025Context} & 
        Medir RHR durante vida diaria: impacto contexto conductual + criterios metodológicos & 
        Estudio observacional & 
        Wearables (RHR + acelerómetro) & 
        --- & 
        Análisis contextual & 
        RHR por contexto & 
        \textbf{RHR sedentario $\neq$ RHR sueño}. Contextos autonómicos diferentes. Clasificación fisiológica vs postural \\
        
        \midrule
        
        Marino et al. (2024) \cite{Marino2024ARIC} & 
        Asociaciones PA + HRV con función cognitiva + demencia (ARIC Neurocognitive Study) & 
        Análisis cohorte longitudinal & 
        ECG ambulatorio 2 semanas + acelerómetro & 
        961 & 
        Análisis regresión & 
        Odds Ratio (OR) & 
        Mayor PA + HRV $\rightarrow$ mejor cognición. Relación no lineal PA-HRV-salud cerebral \\
        
        \midrule
        
        % ========== INVESTIGACIÓN ACTUAL (ÚLTIMA FILA) ==========
        
        \textbf{Martínez-Corral (2025)} \newline \textbf{[Investigación Actual]} & 
        \textbf{Clasificación comportamiento sedentario mediante lógica difusa + datos biométricos Apple Watch (HRV, acelerómetro)} & 
        \textbf{Clustering K-Means ($k=2$) $\rightarrow$ Fuzzy Mamdani (4 variables)} & 
        \textbf{Apple Watch (PPG + IMU) vida libre} & 
        \textbf{10 (5F/5M)} & 
        \textbf{LOUO (10-fold)} & 
        \textbf{F1-Score} & 
        \textbf{F1-global=0.840, F1-LOUO=0.780$\pm$0.167. Interpretabilidad 100\%. BYOD validado. Único estudio clustering$\rightarrow$fuzzy + LOUO + HRV sedentarismo} \\
        
        \bottomrule
    \end{tabular}
\end{table}

% ============================================================================
% NOTAS PARA EL AUTOR:
% ============================================================================
% 
% REFERENCIAS CITADAS EN LA TABLA (17):
% 1. Capitoli2025FuzzyXAI
% 2. MarashiHosseini2023Dietary
% 3. Goncalves2021
% 4. Khan2024Wearable
% 5. Bassani2025DLHAR
% 6. Fuller2021Predicting
% 7. Mathew2024LOSO
% 8. Rehman2024LOSO
% 9. Salim2024STEPHEN
% 10. Migueles2022GRANADA
% 11. Lyons2024StandHour
% 12. OGrady2024AppleWatch
% 13. Ji2023Scratch
% 14. Bienefeld2023XAI
% 15. Deng2023LharJBHI
% 16. Godkin2025Context
% 17. Marino2024ARIC
% 
% ARTÍCULOS QUE USAN LOUO/LOSO (5):
% - Bassani2025DLHAR (LOSO caída rendimiento)
% - Fuller2021Predicting (LOOCV)
% - Mathew2024LOSO (F1 drop personalizado→LOSO)
% - Rehman2024LOSO (comparación LOSO vs k-fold)
% - Ji2023Scratch (LOSO Nature npj)
% 
% DISTRIBUCIÓN METODOLÓGICA:
% - Fuzzy Logic: 3 artículos (Capitoli, Marashi, Gonçalves)
% - Deep Learning: 2 artículos (Khan, Bassani)
% - ML clásico: 3 artículos (Fuller, Mathew, Rehman)
% - Métodos estadísticos: 2 artículos (Salim, Migueles)
% - Validación wearables: 3 artículos (Lyons, OGrady, Ji)
% - XAI: 2 artículos (Bienefeld, Deng)
% - Fisiología/HRV: 2 artículos (Godkin, Marino)
% 
% APORTACIÓN DIFERENCIAL (investigación actual):
% 1. ÚNICO estudio clustering NO SUPERVISADO → fuzzy SUPERVISADO + LOUO
% 2. ÚNICO estudio con HRV para clasificar sedentarismo (fisiológico vs postural)
% 3. Validación LOUO (generalización inter-usuario)
% 4. Interpretabilidad 100% (reglas lingüísticas + funciones membresía)
% 5. BYOD (Bring-Your-Own-Device) vida libre multi-anual
% 
% ============================================================================

